% !TEX root = ../main.tex

\chapter{总结与展望}

本文围绕“融合三视图特征的鞋履3D模型重建研究”这一课题，面向鞋履设计草图到三维网格模型的快速生成需求，设计并实现了一个本地可运行的鞋履三维重建原型系统。系统以单视图鞋履草图到三维GLB模型作为当前稳定闭环，同时在数据组织、方法设计和后续接口中保留侧视图、俯视图和底视图的融合方向。系统结合草图预处理、ControlNet草图域适配、SF3D基线、Hunyuan3D-2mini形状生成、Gradio演示界面和非破坏式网格质量报告，实现了从二维输入到三维预览、指标分析和论文材料生成的完整流程。本章对全文工作进行总结，归纳主要结论，并讨论后续可以继续改进的方向。

\section{工作总结}

本文首先分析了鞋履设计草图到三维模型重建的研究背景和任务特点。鞋履设计草图能够低成本表达鞋面轮廓、鞋底结构、鞋带区域和后跟形态，但二维草图难以直接展示三维体积关系和内部结构。传统三维建模流程需要较多人工建模和修正成本，而近年来图像到三维生成模型和扩散先验方法为快速生成可预览三维资产提供了新的可能。基于这一背景，本文将研究目标定位为在本地硬件条件下构建可运行的鞋履外观级三维网格生成系统，并为三视图融合和结构约束继续扩展打下基础。

其次，本文梳理了单图像三维重建、扩散先验三维生成、多视图生成、前馈式三维重建和草图条件控制等相关工作。在模型路线选择上，本文综合考虑开源可获得性、本地部署复杂度、GLB导出稳定性和鞋履样例视觉效果，采用“SF3D基线+Hunyuan3D-2mini主后端+ControlNet草图域适配”的路线。SF3D用于建立稳定参照和回退路径，Hunyuan3D-2mini用于承担主要形状生成，ControlNet用于缓解草图输入与三维生成模型常见输入之间的域差异。

再次，本文设计了面向鞋履草图输入的端到端方法流程。该流程包括图像白底合成、方形补边、尺寸归一化、灰度化、二值线稿提取、Canny和Scribble控制图生成等预处理步骤；同时引入ControlNet作为草图域适配模块，将线稿式鞋履草图转换为更接近产品渲染图的中间图像。三维重建阶段支持SF3D和Hunyuan3D-2mini两类后端，并通过统一的CLI流水线管理直接输入、ControlNet输入和多后端对比。

在工程实现方面，本文构建了模块化项目结构。系统以WSL作为主要机器学习运行环境，通过主控环境、\texttt{hunyuan3d}和\texttt{sf3d}等conda环境隔离不同模型依赖。核心CLI程序负责连接预处理、ControlNet渲染、Hunyuan3D/SF3D推理和网格报告生成。Gradio界面作为轻量交互层调用同一CLI入口，提供图像上传、参数设置、GLB预览、中间图像展示、网格指标表和警告表等功能。该设计保证了命令行实验和网页演示使用同一套后端逻辑，提升了结果可复现性。

最后，本文设计并实现了非破坏式网格质量报告模块。该模块读取生成的GLB文件，统计几何数量、顶点数、面数、包围盒尺寸、厚度代理比例、连通组件、水密性、表面积和体积等指标，并生成结构化风险提示。特别地，模块始终保留内部结构不可验证警告，用于提醒单视图或有限视图条件下无法仅凭自动指标判断鞋履内部空间是否合理。该报告既服务于Gradio演示，也为论文实验分析提供了可复现的量化依据。

\section{主要结论}

通过系统实现和冻结样例实验，本文得到以下结论。

第一，在本地移动端GPU环境中，构建鞋履草图到三维GLB模型的演示系统是可行的。实验中的Hunyuan3D-2mini shape-only路径能够在约4.4GB峰值CUDA显存下生成可加载GLB模型，Gradio界面验证样例也完成了从输入图像到三维预览和网格报告的完整流程。这说明在合理选择模型后端和推理设置的情况下，本地轻量级三维生成演示能够满足毕业设计原型系统的基本需求。

第二，SF3D适合作为稳定基线，但不适合作为本文的最佳结果路线。SF3D能够较快生成可加载GLB文件，便于验证预处理、ControlNet适配、CLI编排和网格报告模块是否贯通。然而，直接草图输入下的SF3D结果厚度代理比例较低，鞋履内部空间和局部几何表达有限；ControlNet渲染输入能够改善其体积感，但仍主要适合作为基线和回退后端。

第三，Hunyuan3D-2mini更适合作为当前系统主后端。冻结样例中，Hunyuan3D-2mini生成的S1和S2模型具有更高的网格复杂度和更连续的外观形状，其中S1包含96774个顶点和193544个面，S2包含117700个顶点和235396个面。结合人工视觉审阅，Hunyuan3D-2mini在当前鞋履样例上的外观表现明显优于SF3D基线，因此更适合作为论文和演示的主路线。

第四，ControlNet草图域适配具有必要性，但需要谨慎使用。SF3D对直接草图输入的结果较薄，S3的厚度代理比例约为0.151731；当输入换成ControlNet渲染图后，S4的厚度代理比例提升到约0.348921，说明渲染图能够提供更多体积线索，改善直接线稿输入导致的扁平化问题。另一方面，ControlNet生成图可能引入原始草图中不存在的细节，因此在鞋履设计任务中不能只追求渲染图真实感，还应关注设计轮廓和结构线的保持程度。

第五，三视图融合和内部结构表达仍是后续完善重点。本文系统已经保留侧视图、俯视图和底视图的数据组织方式，并在方法设计中讨论了多视图约束的价值。当前稳定闭环主要验证单视图输入下的本地生成流程，后续需要进一步将三视图信息接入后端输入、融合策略和一致性评价中，以更充分地约束鞋头、鞋底、鞋口和内部空腔等结构。

\section{未来工作}

虽然本文已经完成了本地可运行的鞋履三维重建原型系统，但距离稳定、真实和可制造的鞋履三维重建仍有较大提升空间。后续工作可以从以下几个方向展开。

第一，完善鞋履草图与三维模型配对数据。后续可以继续整理鞋履设计草图、三视图图纸、商品渲染图和对应三维模型，并区分侧视图、俯视图、底视图和细节图。若能够形成较稳定的数据集，将可以支持更严格的监督训练、模型微调和定量评价。

第二，推进ControlNet鞋履草图域适配。当前ControlNet使用通用Canny控制和通用提示词，能够生成较可用的中间渲染图，但仍可能改变原始设计细节。后续可以使用鞋履草图-渲染图数据对ControlNet进行小规模微调，使其更好地保持鞋履轮廓、鞋底层次、鞋带区域和设计线条。同时，可以比较Canny、Scribble、深度、法线等不同控制形式对鞋履重建结果的影响。

第三，继续探索Hunyuan3D多视图和微调路线。后续可以进一步研究如何将侧视图、俯视图和底视图共同作为输入约束三维形状，并在更充分的数据和实验条件下评估多视图生成对鞋口、鞋底厚度和内部空腔的改善作用。若具备足够数据和计算资源，还可以探索Hunyuan3D相关模型的鞋履领域微调，使模型更熟悉鞋履类别的比例、厚度和结构特点。

第四，引入内部结构和厚度约束。鞋履模型不能只关注外部轮廓，还需要合理的鞋口、鞋腔和鞋帮厚度。未来可以尝试引入鞋楦或足部模型作为内部空间先验，或在生成后处理中加入最小厚度、鞋口开口、内外表面分离等结构约束。对于三视图输入，也可以通过侧视图和俯视图联合约束鞋宽与鞋腔范围，减少隐藏几何的随意补全。

第五，改进网格后处理和模型可用性。当前系统默认只生成非破坏式网格报告，不直接修改几何。后续可以在保留原始GLB的基础上加入可选的平滑、减面、碎片清理、法线修复和尺度归一化模块，并通过人工审阅和指标对比判断这些操作是否改善模型质量。若目标进一步接近设计生产流程，还需要研究从生成网格到参数化曲面或CAD草模的转换方法。

第六，完善评价体系和用户验证。本文实验样例数量有限，主要用于验证系统闭环和初步比较后端表现。后续可以扩展更多鞋型、视角和草图风格，设计更系统的评价指标，例如轮廓保持度、多视图一致性、网格完整性、人工视觉评分和设计师可用性问卷等。若能获得真实三维模型，还可以加入Chamfer Distance、法向一致性和体积误差等几何指标。

综上，本文完成了从鞋履设计草图到三维GLB模型的本地原型系统，并通过实验验证了SF3D、Hunyuan3D-2mini和ControlNet组合路线的可行性与边界。未来若能结合鞋履领域数据、多视图约束、结构先验和模型微调，该方向有望进一步提升从设计草图到三维鞋履资产生成的质量和可控性。
